% Entregable 4 — documento técnico de UNA página (vertical). Compilar en Overleaf (pdflatex).
% Todas las cifras salen de results/ vía `bash scripts/run_e4.sh tablas` y de ft/lector/ (ver docs/e4_finetuning.md).
% En rojo (\pend{}): lo que el equipo debe completar antes de entregar (video, Hugging Face, revisión a mano).
% Si no cabe en una página: bajar `height` del gráfico de pérdida o quitar la última frase de la sección 6.
% En \texttt, un guion doble se escribe -{}- (si no, sale como raya).
\documentclass[9pt,a4paper]{extarticle}
\usepackage[T1]{fontenc}
\usepackage[utf8]{inputenc}
\usepackage{lmodern}
\usepackage[spanish,es-noshorthands,es-tabla]{babel}
\usepackage[a4paper,margin=1.0cm]{geometry}
\usepackage{microtype}
\usepackage{booktabs}
\usepackage{tabularx}
\usepackage{multicol}
\usepackage{enumitem}
\usepackage{xcolor}
\usepackage{tikz}
\usetikzlibrary{positioning,arrows.meta,calc,fit}
\usepackage{pgfplots}
\pgfplotsset{compat=1.15}
\usepackage[hidelinks]{hyperref}
\usepackage{titlesec}

\definecolor{sec}{RGB}{30,50,170}
\titleformat{\section}{\normalsize\bfseries\color{sec}}{\thesection.}{0.4em}{}
\titlespacing*{\section}{0pt}{0.25em}{0.1em}
\setlist{nosep,leftmargin=1.1em}
\setlength{\parindent}{0pt}
\setlength{\parskip}{0.15em}
\setlength{\tabcolsep}{3pt}
\setlength{\columnsep}{0.5cm}
\pagestyle{empty}

\newcommand{\repo}{\url{https://github.com/benjamin-cpu-alt/Generative-AI-UdeC}}
\newcommand{\video}{\textcolor{red}{\textbf{[LINK DEL VIDEO]}}}
% Reemplazar USUARIO_HF por la cuenta donde se publicó el adaptador (scripts/ft/subir_hf.py).
\newcommand{\hf}{\url{https://huggingface.co/USUARIO_HF/phi4-mini-lector-lora}}
\newcommand{\pend}[1]{\textcolor{red}{#1}}
\newcommand{\cod}[1]{\texttt{#1}}
\newcommand{\warn}{\textcolor{orange!80!black}{\textbf{(!)}}}
\newcommand{\alto}{\textcolor{red!70!black}{\rule{0.9ex}{0.9ex}}}% condición de parada

\begin{document}
\begin{center}
{\large\bfseries Entregable 4: Fine-tuning del agente Lector (LoRA sobre Phi-4-mini)}\\[1pt]
\textbf{Equipo:} Benjamín Grandón V., Carla Maureira V. — Universidad de Concepción, Generative AI (580694)\\
\textbf{Repositorio:} \repo \quad \textbf{Video:} \video\\
\textbf{Pesos (adaptador):} \hf
\end{center}
\vspace{-5pt}

% ------------------------------------------------------------------ grafo --
\begin{center}
\begin{tikzpicture}[
  font=\scriptsize, node distance=5mm and 11mm,
  ag/.style={draw, rounded corners=3pt, fill=blue!8, draw=blue!60!black, align=center, inner sep=3pt, text width=27mm, minimum height=13mm},
  ft/.style={ag, fill=orange!22, draw=orange!85!black, line width=1.3pt},
  tool/.style={draw, fill=green!9, draw=green!45!black, align=center, inner sep=2pt, text width=21mm, minimum height=8mm},
  io/.style={draw, fill=gray!8, draw=gray!60!black, align=center, inner sep=2pt, text width=19mm, minimum height=8mm},
  stop/.style={draw, double, fill=red!6, draw=red!60!black, align=center, inner sep=2pt, text width=16mm},
  arr/.style={-{Stealth[length=1.8mm]}, thick},
  loop/.style={-{Stealth[length=1.8mm]}, thick, dashed, red!60!black},
  lab/.style={font=\tiny, align=center, fill=white, inner sep=1pt}]
\node[io] (in) {Aviso (texto crudo), uno a la vez};
\node[ft, right=of in] (lec) {\textbf{A1 · Lector \,$\star$ LoRA}\\phi4-mini \textbf{+ adaptador}: copia el fragmento literal de cada tema};
\node[tool, right=of lec] (anc) {\cod{anclar}\\fragmentos $\to$ datos; marca \textbf{dudas} (\warn{})};
\node[ag, right=of anc] (ver) {\textbf{A2 · Verificador} (base)\\por campo \warn{}: confirmar · corregir · releer · no\_verificable};
\node[tool, below=9mm of ver] (com) {\cod{comprobar}\\¿la evidencia está en el aviso y da ese valor?};
\node[tool, below=9mm of anc] (eva) {\cod{evaluar}\\5 restricciones por aviso};
\node[ag, below=9mm of lec] (dec) {\textbf{A3 · Decisor} (base)\\por caso: cerrar o consultar un dato};
\node[tool, left=of dec] (ord) {\cod{ordenar}\\JSON E1 (preferida, ROI)};
\node[stop, below=5mm of ord] (out) {\alto{} \textbf{FIN}: JSON estricto};
\node[io, below=3mm of in] (perfil) {Perfil del comprador};

\draw[arr] (in) -- node[lab, above=1pt]{texto} (lec);
\draw[arr] (lec) -- node[lab, above=1pt]{spans} (anc);
\draw[arr] (anc) -- node[lab, above=1pt]{hechos +\\\warn{} y frase} (ver);
\draw[arr] (ver) -- node[lab, right]{campo, evidencia,\\valor} (com);
\draw[loop] (com.east) to[out=0,in=0,looseness=1.6] node[lab, right]{\textbf{reintento}:\\rechazo + motivo\\($\le$1; luego \alto{} null)} (ver.east);
\draw[loop] (ver.north) to[out=150,in=30] node[lab, above]{\textbf{relectura}: campos + motivo ($\le$2 vueltas; luego \alto{} dato null)} (lec.north);
\draw[arr] (com.west) -- node[lab, above]{\alto{} hechos\\finales} (eva.east);
\draw[arr] (anc.south) -- node[lab, right]{\alto{} sin \warn{}:\\hechos} (eva.north);
\draw[arr] (eva) -- node[lab, above]{estado por\\restricción} (dec);
\draw[loop] (dec.south) to[out=-30,in=-120] node[lab, below]{\textbf{consulta}(id, campo) si no\_verificable o sin\_confirmar ($\le$1 ronda; luego \alto{} cierre forzado)} (ver.south west);
\draw[arr] (dec) -- node[lab, above]{cerrar} node[lab, below]{\alto{}} (ord);
\draw[arr] (ord) -- node[lab, right=1pt]{JSON E1} (out);
\draw[arr, gray] (perfil.east) -| node[lab, pos=0.3]{restricciones\\+ preferencias} (dec.north);
\end{tikzpicture}
\end{center}
\vspace{-7pt}
{\scriptsize \textbf{Naranja ($\star$): el agente ajustado}; azul: agentes con el phi4-mini base; verde: herramientas Python, sin modelo; rojo discontinuo: los tres bucles y su máximo; \alto{}: parada.}

\begin{multicols}{2}
\section{Continuidad y elección del agente}
Misma tarea, juez y entradas que E1--E3. \textbf{¿Qué agente?} Re-ejecutamos el orquestador con las salidas grabadas en las trazas de la E3 sobre \cod{hard\_test} (\cod{scripts/atribucion\_e3.py}): de 218 errores de primera lectura, \textbf{162 nacen en el Lector} (117 son dormitorios en \cod{null} ante «3D», «2 hab.») y 56 en la herramienta \cod{anclar}. \textbf{¿Por qué pesos y no más prompt?} El prompt ya pide copiar ese fragmento: no falta una regla, el modelo no reconoce esas formas; y en la E2 iterar el prompt empeoró (21 $\to$ 13 $\to$ 8 $\to$ 8/30 en dev). El \textbf{control} lo mide: Lector base + 8 ejemplos few-shot de los mismos datos.

\section{Cómo corren los bucles (comentario de la E3)}
\textbf{Agente} = llamada al modelo con su prompt; devuelve una \emph{acción} de un \cod{enum} que decide a quién va el trabajo. \textbf{Herramienta} = función Python: calcula o valida y devuelve su resultado; nunca elige el paso siguiente. El \textbf{orquestador} es un \cod{while} de Python, no un agente: lee la acción, entrega el mensaje y cuenta vueltas. \textbf{Relectura}: el Verificador responde \cod{releer} y el Lector relee esos campos ($\le$2). \textbf{Reintento}: \cod{comprobar} rechaza con un motivo y el Verificador decide otra vez ($\le$1). \textbf{Consulta}: el Decisor responde \cod{consultar(id, campo)} y el Verificador revisa ese dato ($\le$1). \textbf{Parada}: \cod{cerrar}, o presupuesto agotado. \cod{agents.show -{}-roles} lo marca en cada traza y en el video.

\section{Datos de entrenamiento}
\textbf{1.146 pares} entrada$\to$salida del rol del Lector (1.038 entrenamiento, 108 validación). \textbf{1.083 sintéticos} (\cod{scripts/ft/build\_lector\_data.py}): frases propias en 6 estilos con el fragmento a copiar marcado, así que la salida es correcta \emph{por construcción}. \textbf{63 autoetiquetados}: salidas del phi4-mini base sobre \cod{train} que coinciden con el \cod{truth}. Ningún modelo más grande generó ni etiquetó. \textbf{Revisión a mano}: 60 al azar, \pend{XX/60 correctos}. \textbf{Separación}: se descarta todo ejemplo con 8 palabras seguidas en común con un set de evaluación (4, todos de \cod{hard\_dev}); un test lo repite.

\section{Entrenamiento}
{\scriptsize
\begin{minipage}[c]{0.47\columnwidth}
\resizebox{\linewidth}{!}{%
\begin{tabular}{@{}ll@{}}
\toprule
Método & LoRA, \cod{mlx-lm} 0.29.1\\
Módulos & atención + MLP, 16 capas\\
Rango / $\alpha$ / drop. & 16 / 32 / 0,05\\
Parámetros & 11,5 M (0,30\,\%)\\
Optimizador & AdamW, wd 0,01\\
Tasa & 1,5e-4, warmup + coseno\\
Lote efectivo & 4 (1 $\times$ acum. 4)\\
Ejemplos vistos & 1.300 (1,25 épocas)\\
Hardware & MacBook Air M4, 16 GB\\
Tiempo / mem. & \textbf{202 min} / 6,7 GB\\
\bottomrule
\end{tabular}}
\end{minipage}\hfill
\begin{minipage}[c]{0.50\columnwidth}
\begin{tikzpicture}
\begin{axis}[width=\linewidth, height=3.2cm, ymode=log, xmin=0, xmax=1300, ymin=0.0015, ymax=1,
  xlabel={ejemplos vistos (pérdida en escala log)}, tick label style={font=\tiny}, label style={font=\tiny},
  xlabel style={yshift=4pt}, grid=major, grid style={gray!20},
  legend style={font=\tiny, draw=none, fill=none, at={(0.98,0.97)}, anchor=north east}, legend cell align=left]
\addplot[blue!55, thin] coordinates {(20,0.473)(40,0.233)(60,0.112)(80,0.076)(100,0.087)(120,0.056)(140,0.083)(160,0.062)(180,0.046)(200,0.073)(220,0.05)(240,0.034)(260,0.071)(280,0.064)(300,0.02)(320,0.03)(340,0.043)(360,0.016)(380,0.012)(400,0.018)(420,0.023)(440,0.032)(460,0.008)(480,0.035)(500,0.029)(520,0.012)(540,0.02)(560,0.014)(580,0.019)(600,0.011)(620,0.019)(640,0.008)(660,0.031)(680,0.03)(700,0.018)(720,0.04)(740,0.018)(760,0.011)(780,0.011)(800,0.01)(820,0.016)(840,0.005)(860,0.007)(880,0.005)(900,0.013)(920,0.016)(940,0.018)(960,0.016)(980,0.004)(1000,0.01)(1020,0.004)(1040,0.01)(1060,0.006)(1080,0.003)(1100,0.004)(1120,0.005)(1140,0.006)(1160,0.003)(1180,0.003)(1200,0.005)(1220,0.002)(1240,0.004)(1260,0.003)(1280,0.002)(1300,0.004)};
\addplot[red!75!black, thick, mark=*, mark size=1.1pt] coordinates {(1,0.69)(260,0.028)(520,0.017)(780,0.01)(1040,0.007)(1300,0.006)};
\legend{entrenamiento, validación (108)}
\end{axis}
\end{tikzpicture}
\end{minipage}}

Pérdida solo sobre la respuesta; se usa el último checkpoint. El adaptador (GGUF, 23 MB) se monta en Ollama sobre el mismo \cod{phi4-mini:latest} de la E3.

\section{Resultados (mismas entradas, mismo juez)}
{\scriptsize
\begin{tabularx}{\columnwidth}{@{}Xccccc@{}}
\toprule
\cod{hard\_test} (30 casos, 240 avisos) & E1 & full & aviso & aprob. ind. & rech. falso\\
\midrule
Rechazar todo (piso) & 30 & 0 & 127 & 0/127 & 113/113\\
E1 · baseline & 1 & 0 & 118 & 38/127 & 84/113\\
E2 · extractor + herramientas & 11 & 1 & 193 & 0/127 & 47/113\\
E3 · agentes (Lector base) & \textbf{23} & 5 & 214 & \textbf{6/127} & 20/113\\
\textbf{E4 · Lector con LoRA} & 22 & \textbf{16} & \textbf{225} & 8/127 & \textbf{7/113}\\
Control · base + 8 few-shot & \textbf{23} & 11 & 221 & 7/127 & 12/113\\
\bottomrule
\end{tabularx}

\begin{tabularx}{\columnwidth}{@{}Xccc@{}}
\toprule
Dentro del sistema (\cod{hard\_test}) & E3 & \textbf{E4} & Control\\
\midrule
Errores de lectura con origen en el Lector & 162 & \textbf{22} & 26\\
Campos que llegan mal a \cod{evaluar} & 75 & \textbf{39} & 47\\
Llamadas al Verificador / reintentos & 315 / 93 & \textbf{134 / 36} & 151 / 42\\
Tokens de entrada / minutos & 577k / 66 & \textbf{368k / 43} & 731k / --\\
\bottomrule
\end{tabularx}

\begin{tabularx}{\columnwidth}{@{}Xcccc@{}}
\toprule
Regresión: full · aprob. ind. & E2 & E3 & \textbf{E4} & Control\\
\midrule
\cod{real\_e3} (20 casos, avisos reales) & 18 · 0 & \textbf{20 · 0} & 19 · 1 & 19 · 1\\
\cod{ood} (12 casos) & 7 · 0 & 10 · 0 & \textbf{11 · 0} & 10 · 0\\
\cod{test} (51 casos) & 47 · 0 & 50 · 0 & \textbf{51 · 0} & \textbf{51 · 0}\\
\bottomrule
\end{tabularx}}

\textbf{Lectura.} Los errores del Lector bajan de 162 a 22 y las respuestas exactas suben de 5 a 16/30 (11 casos a favor, 0 en contra; McNemar $p{=}0{,}001$). \textbf{Lo que los pesos consiguen y el prompt no} es más chico: el control ya baja esos errores a 26; E4 le gana por 5 casos exactos (5 a 0, $p{=}0{,}06$) y 5 rechazos falsos, con la mitad de tokens: consistente, pero no concluyente con 30 casos. Costo: 2 aprobaciones indebidas nuevas en \cod{hard\_test} y 1 en los avisos reales; \cod{test} y \cod{ood} no empeoran.

\section{Caso de fallo}
\cod{HT-0906} (\cod{results/agents\_ft\_hard\_test/ht\_09}): «Macul, 2D. [\dots] Estacionamiento: no; hay estacionamiento en calle.»; el comprador exige estacionamiento. Con el Lector base pasaba por el Verificador y se rechazaba bien. El Lector ajustado \textbf{sí copia «2D»}, pero corta en «Estacionamiento: no»; \cod{anclar} lo lee como \emph{propio} sin marcar duda, el aviso \textbf{se salta al Verificador} y se \textbf{aprueba indebidamente}. \textbf{Aviso real} \cod{REAL-38}: el Lector ajustado lee mejor la comuna, el Verificador base recibe otro mensaje, inventa «propio» citando «estación de metro» y \cod{comprobar} lo acepta. Ajustar un agente cambia lo que ven los demás; no corregimos nada tras ver los sets de evaluación.

\section{Reproducibilidad}
README, sección E4: \cod{bash docs/demo\_video\_e4.sh} reproduce el video y \cod{bash scripts/run\_e4.sh tablas} regenera las tablas sin modelo. \cod{SHA256SUMS} fija los pesos evaluados.
\end{multicols}
\end{document}
